% Sección 2 (tarea C1) - autor: Persona C. Editar solo este archivo.
\section{Especificación del lenguaje}

\subsection{Propósito}
DomoDSL permite escribir programas formados por una o más \textbf{reglas de automatización}. Cada regla evalúa una
\textbf{condición} sobre variables del hogar (temperatura, humedad, movimiento, hora, \dots) y, si se cumple, ordena
una o más \textbf{acciones} sobre dispositivos (ventilador, alarma, termostato, \dots). El lenguaje no ejecuta las
reglas: el objetivo del proyecto es \emph{reconocerlas} (decidir si son sintácticamente válidas), informar errores
útiles y producir una representación estructurada (AST).

\subsection{Alfabeto}
El alfabeto relevante $\Sigma$ está formado por:
\begin{itemize}[nosep]
\item letras \texttt{A--Z} y \texttt{a--z}, dígitos \texttt{0--9}, el guion bajo \texttt{\_} y el punto \texttt{.} (solo dentro de números decimales);
\item los símbolos \texttt{>}, \texttt{<}, \texttt{=}, \texttt{!}, \texttt{;}, \texttt{,}, \texttt{(} y \texttt{)};
\item espacio, tabulación y salto de línea (separadores que el lexer descarta).
\end{itemize}
Cualquier otro carácter produce un \emph{error léxico}.

\subsection{Categorías de tokens}
El analizador léxico (basado en el entregado por el docente) usa expresiones regulares. Los espacios se descartan.
Las palabras reservadas \textbf{no distinguen mayúsculas de minúsculas} (como en la función base).

\begin{center}\small
\begin{tabular}{@{}>{\raggedright\arraybackslash}p{2.5cm} >{\raggedright\arraybackslash}p{4.0cm} >{\raggedright\arraybackslash}p{4.3cm} >{\raggedright\arraybackslash}p{3.3cm}@{}}
\toprule
\textbf{Token} & \textbf{Patrón} & \textbf{Descripción} & \textbf{Ejemplos} \\
\midrule
\texttt{WHEN, THEN} & \texttt{WHEN | THEN} & Inicio y separador de la regla & \texttt{WHEN, then} \\
\texttt{AND, OR, NOT} & \texttt{AND | OR | NOT} & Operadores lógicos & \texttt{AND, not} \\
\texttt{TRUE, FALSE} & \texttt{TRUE | FALSE} & Valores booleanos & \texttt{TRUE} \\
\texttt{ID} & \texttt{[A-Za-z\_][A-Za-z0-9\_]*} & Identificadores (variables, verbos y dispositivos) & \texttt{temperatura, sensor\_1} \\
\texttt{NUM} & \texttt{\textbackslash{}d+(\textbackslash{}.\textbackslash{}d+)?} & Números enteros o decimales & \texttt{28, 0.5} \\
\texttt{OP\_REL} & \texttt{>= | <= | == | != | > | <} & Operadores relacionales & \texttt{>=, ==} \\
\texttt{PUNTO\_COMA} & \texttt{;} & Fin de regla & \texttt{;} \\
\texttt{COMA} & \texttt{,} & Separador de acciones & \texttt{,} \\
\texttt{PARENT\_IZQ} & \texttt{(} & Abre agrupación & \texttt{(} \\
\texttt{PARENT\_DER} & \texttt{)} & Cierra agrupación & \texttt{)} \\
\bottomrule
\end{tabular}
\end{center}

Un identificador cuyo texto coincide con una palabra reservada se convierte en el token reservado
(p.\,ej. \texttt{when} $\to$ \texttt{WHEN}), por lo que las palabras reservadas no pueden usarse como identificadores.
Al final de la entrada el lexer agrega el token \texttt{\$} (fin de entrada), que necesita el analizador LL(1).

\subsection{Cumplimiento de los requisitos mínimos del enunciado}
\begin{center}\small
\begin{tabular}{@{}>{\raggedright\arraybackslash}p{5.6cm} >{\raggedright\arraybackslash}p{9.4cm}@{}}
\toprule
\textbf{Requisito} & \textbf{Cómo lo cumple DomoDSL} \\
\midrule
Al menos 4 palabras reservadas & 7: \texttt{WHEN}, \texttt{THEN}, \texttt{AND}, \texttt{OR}, \texttt{NOT}, \texttt{TRUE}, \texttt{FALSE} \\
Identificadores & Token \texttt{ID} \\
Valores numéricos y/o booleanos & \texttt{NUM} (enteros y decimales) y \texttt{TRUE}/\texttt{FALSE} \\
Al menos 3 operadores & 9: seis relacionales (\texttt{> < >= <= == !=}) y tres lógicos (\texttt{AND OR NOT}) \\
Estructura de sentencia principal & \texttt{WHEN <condición> THEN <acciones> ;} \\
Más de una sentencia & Un programa es una secuencia de una o más reglas \\
Construcción compuesta & Conectores \texttt{AND}/\texttt{OR}, paréntesis y listas de acciones con parámetros \\
Extensión creativa & \texttt{NOT}, paréntesis, múltiples acciones y acciones con parámetros (sección~4.6) \\
\bottomrule
\end{tabular}
\end{center}

\subsection{Ejemplos de sentencias válidas e inválidas}
\paragraph{Válidas.}
\begin{enumerate}[nosep]
\item \texttt{WHEN temperatura > 28 THEN encender ventilador;}
\item \texttt{WHEN movimiento == TRUE THEN activar alarma;}
\item \texttt{WHEN temperatura > 28 AND humedad < 40 THEN encender ventilador;}
\item \texttt{WHEN luz < 100 OR hora >= 20 THEN encender lampara;}
\item \texttt{WHEN NOT puerta == TRUE THEN activar alarma;}
\item \texttt{WHEN (temperatura > 30 OR humedad > 80) AND presencia == TRUE THEN encender aire;}
\item \texttt{WHEN temperatura > 28 THEN encender ventilador 3, activar alarma;}
\end{enumerate}
\paragraph{Inválidas.}
\begin{center}\small
\begin{tabular}{@{}l >{\raggedright\arraybackslash}p{7.2cm} >{\raggedright\arraybackslash}p{6.3cm}@{}}
\toprule
\textbf{N.º} & \textbf{Sentencia} & \textbf{Motivo} \\
\midrule
1 & \texttt{WHEN temperatura > 28 encender ventilador;} & Falta la palabra reservada \texttt{THEN} \\
2 & \texttt{WHEN temperatura > THEN encender ventilador;} & Falta el valor después de \texttt{>} \\
3 & \texttt{temperatura > 28 THEN encender ventilador;} & La regla no comienza con \texttt{WHEN} \\
4 & \texttt{WHEN temperatura > 28 THEN encender ventilador} & Falta el \texttt{;} final \\
5 & \texttt{WHEN (temperatura > 28 THEN encender ventilador;} & Paréntesis sin cerrar \\
6 & \texttt{WHEN temperatura = 28 THEN encender ventilador;} & \texttt{=} no es un operador (debe ser \texttt{==}): error léxico \\
\bottomrule
\end{tabular}
\end{center}

